\section{Feedback Design and Supervisory Constraints}
\label{sec:control}
\subsection{Discrete Integral State Feedback}
\textbf{Linear quadratic integral (LQI) control} augments the local state with $\dot z=\rho-d^\ast$. For $\boldsymbol\xi=[\mathbf x^\top,z]^\top$,
\begin{equation}
A_a=\begin{bmatrix}A&0\\C_r&0\end{bmatrix},\quad
B_a=\begin{bmatrix}B\\0\end{bmatrix},\quad C_r=\begin{bmatrix}1&0&0\end{bmatrix}.
\end{equation}
Exact zero-order-hold discretization at $T_s$ uses a block exponential,
\begin{equation}
\exp\!\left(T_s\begin{bmatrix}A_a&B_a\\0&0\end{bmatrix}\right)
=\begin{bmatrix}A_d&B_d\\0&I_2\end{bmatrix}.
\end{equation}
This avoids inverting the singular matrix $A_a$. The nominal infinite-horizon cost is
\begin{equation}
J_c=\sum_{k=0}^{\infty}\left(\boldsymbol\xi_k^\top Q\boldsymbol\xi_k+\mathbf u_k^\top R\mathbf u_k\right).
\end{equation}
The diagonal weights are $Q=\operatorname{diag}(6,2,0.8,0.3)$ and $R=\operatorname{diag}(25,12)$ in the stated SI coordinate convention; their units are reciprocal to the squared state/input units, so each cost term is dimensionless. The integral state has units $\mathrm{m\,s}$. The Riccati equation and gain are
\begin{equation}
\begin{aligned}
P_R&=A_d^\top P_RA_d+Q\\
&\quad-A_d^\top P_RB_d(R+B_d^\top P_RB_d)^{-1}B_d^\top P_RA_d,\\
K_c&=(R+B_d^\top P_RB_d)^{-1}B_d^\top P_RA_d.
\end{aligned}
\end{equation}
The virtual command is $\mathbf u_k=-K_c[\hat\rho-d^\ast,\hat\beta,\hat\varepsilon,z]^\top$. For Table~\ref{tab:assumptions}, the numerical gain is
\begin{equation}
K_c\simeq\begin{bmatrix}
-0.668891&0&0&-0.107705\\
0&-0.402206&-0.186979&0
\end{bmatrix}.
\end{equation}
The negative range coefficients are consistent with $\dot{\delta\rho}=-\delta v_f+\delta v_t$: a positive range error must increase forward speed. Gains are simulation values pending identification.

For the ideal unconstrained full-state discrete loop $F=A_d-B_dK_c$, the poles are approximately $0.977032$, $0.979657$, $0.986764$, and $0.993433$. A discrete Lyapunov solution satisfies
\begin{equation}
F^\top P_LF-P_L=-I_4,\qquad P_L\succ0,
\end{equation}
with minimum eigenvalue $16.3695$ and maximum absolute residual $1.42\times10^{-14}$ in these coordinates. Thus this local reference is asymptotically stable. The certificate excludes saturation, asynchronous estimation, actuator lag, and supervisory switching.

The design discretizes the ideal augmented dynamics exactly. Runtime integration uses a forward-Euler update with anti-windup:
\begin{equation}
z_{k+1}=\sat_{[-z_{\max},z_{\max}]}\left\{z_k+T_s\left[e_k+\frac{v_k-v_{\mathrm{des},k}}{k_iT_{\mathrm{aw}}}\right]\right\},
\end{equation}
where $e_k=\hat\rho_k-d^\ast$, $k_i=-K_{c,14}$, $T_{\mathrm{aw}}=1\,\mathrm s$, $z_{\max}=3\,\mathrm{m\,s}$, and $v_k$ is the realizable forward command. Consequently, the ideal poles are not the eigenvalues of every runtime operating mode. The simulations evaluate the implemented update.

\subsection{Comparisons and Output Feedback}
The \textbf{pole-placement controller} assigns the range/integral pair to $\exp\{[-0.30,-0.45]T_s\}$ and the bearing/heading pair to $\exp\{[-0.35,-0.50]T_s\}$, with continuous pole parameters in $\mathrm{s^{-1}}$. Channel-wise assignment makes one explicit choice among the nonunique MIMO feedback gains. Alternative gains with the same eigenvalues can change eigenvectors, input effort, coupling, and sensitivity.

The \textbf{bearing-only steering baseline} commands $v_{\mathrm{des}}=0.22\,\mathrm{m/s}$ and $\omega_{\mathrm{des}}=0.6\hat\beta\,\mathrm{s^{-1}}$. It shares sensing, estimation, speed limits, and the range-dependent supervisor with the other methods. Its tracking law has no range regulator; the complete baseline is therefore not a purely bearing-only sensing system. Shared stopping logic prevents an unfair comparison against an unconstrained constant-speed vehicle.

In the ideal synchronous linear output-feedback reference, coordinates $[\boldsymbol\xi^\top,\tilde{\mathbf x}^\top]^\top$ yield a block upper-triangular closed-loop matrix with diagonal blocks $F$ and $F_e=(I-L_oC)A_{d,3}$. Here $A_{d,3}$ is the discretized three-state plant, $L_o$ a linear correction gain, and $\tilde{\mathbf x}$ the estimation error. The eigenvalues are the union of controller and observer-error poles. This separation-principle check does not establish stability for the delayed nonlinear EKF and switched constrained controller.

\subsection{Feasible Commands and Terminal Approach}
Forward speed is clipped to the configured maximum and braking bound, reduced by $\max\{0.25,\cos^2\hat\beta\}$, and set to zero below the motor deadband. The yaw command obeys
\begin{equation}
|\omega_f|\le\frac{v_f}{L}\tan\delta_{s,\max},\qquad
\delta_s=\arctan\!\left(\frac{L\omega_f}{v_f}\right),\quad v_f>0.
\end{equation}
At zero speed, both commanded yaw and steering are set to zero. Reverse motion is excluded.

Following and stationary-source approach are distinct missions. The operator must explicitly select the stationary mission, which initializes a zero target-velocity prior and uses $v_{\mathrm{des}}=0.8(\hat\rho-d^\ast)\,\mathrm{s^{-1}}$ and $\omega_{\mathrm{des}}=1.0\hat\beta\,\mathrm{s^{-1}}$. A terminal band accounts for braking distance and motor deadband, then latches the stop. This shared approach mode is used under every controller label; its results cannot demonstrate an LQI advantage at a stationary equilibrium.

\subsection{Supervision and Conditional Stopping Bound}
In ACQUIRE, the chassis remains stationary until three accepted observations pass the estimator and acquisition checks. Rejected packets do not count, but do not necessarily reset the accepted-pair counter. TRACK applies feedback. Stale measurements, excessive covariance, or a bearing outside $\pm55^\circ$ produce recoverable HOLD. A near-target mean range or the explicit terminal condition produces latched STOP. Horn scanning enables reacquisition while the vehicle is stopped.

The speed limiter imposes
\begin{equation}
v_fT+\frac{v_f^2}{2a_b}\le
\hat\rho-3\sigma_\rho-r_{\mathrm{clear}}-m,
\label{eq:brake}
\end{equation}
where $T$ is the loss timeout plus speed lag, $a_b>0$ the assumed braking deceleration, $\sigma_\rho$ the estimated range standard deviation, $r_{\mathrm{clear}}$ the reference-point clearance, and $m$ an additional margin. Every term has units of distance. A negative right-hand side blocks motion. The bound presumes a meaningful uncertainty interval, a valid braking model, and a nonapproaching target. It cannot prevent a target from driving into the follower or detect an obstacle outside the acoustic measurement.

A separate motor process stops when the controller heartbeat expires. This covers a stalled control process while the motor process and driver remain responsive; it does not cover a frozen operating system, blocked motor bus, or failed driver. Physical trials require a separate motor power cut.

\pending{H4 --- Physical constraints and watchdog}{Measure stopping distance at several speeds, heartbeat response, driver stop behavior, steering/speed saturation, and recovery after missing packets. State the measured domain in which the stopping bound is credible.}
